% !TEX root = ../../ctfp-print.tex

\lettrine[lhang=0.17]{是}{时候了}我们该谈谈集合论了。
数学家对集合论有着爱恨交织的关系。它曾是数学的汇编语言——至少过去如此。
范畴论在某种程度上试图远离集合论。例如众所周知所有集合构成的"集合"不存在，
但所有集合构成的范畴 $\Set$ 却存在。这是个好消息。另一方面我们假设任一范畴中任意两个对象间的态射构成一个集合。
我们甚至称其为同态集合(hom-set)。公平地说，确实存在一个范畴论分支研究态射不构成集合的情况：
此时它们是另一个范畴中的对象。这种使用同态对象(hom-objects)而非同态集合的范畴称为\newterm{丰富范畴(enriched categories)}。
不过在下文中我们将坚持使用传统同态集合的范畴。

集合是你能在范畴对象之外找到的最接近无特征团块的东西。集合具有元素，但你无法对这些元素说太多。
若你面对有限集合，你可以数清元素数量。你也可以通过基数来"数"无限集合的元素。
例如自然数集合比实数集合小，尽管两者都是无限的。但或许令人惊讶的是有理数集合与自然数集合大小相同。

除此之外，关于集合的所有信息都可以通过它们之间的函数——特别是可逆函数即同构(isomorphism)来编码。
从所有实际目的来看，同构的集合是相同的。在我招致基础数学家的愤怒之前，请允许我解释：
相等性与同构之间的区别具有根本重要性。事实上这是最新数学分支同伦类型理论(Homotopy Type Theory, HoTT)的主要关注点之一。
我提及HoTT是因为这是一个受计算启发的纯数学理论，其主要倡导者Vladimir Voevodsky在研究Coq定理证明器时获得了重大顿悟。
数学与编程之间的互动是双向的。

关于集合的重要教训是：比较不同元素类型的集合是可以接受的。例如我们可以断言某个自然变换集合与某个态射集合同构，
因为集合就是集合。在这种情况下同构意味着每有一个集合中的自然变换，另一集合中就有唯一的态射与之对应，反之亦然。
它们可以相互配对。当苹果与橙子来自不同范畴时你无法比较苹果与橙子，但你可以比较苹果集合与橙子集合。
经常将范畴问题转化为集合论问题能给我们带来必要洞察力，甚至让我们证明有价值的定理。

\section{Hom函子}

每个范畴都天然配备了一个指向 $\Set$ 的映射族。这些映射实际上是函子，因此它们保持范畴的结构。让我们构建这样一个映射。

设 $\cat{C}$ 中固定对象 $a$，另取对象 $x$ 也在 $\cat{C}$ 中。同态集合 $\cat{C}(a, x)$ 是 $\Set$ 中的对象。
当我们保持 $a$ 固定时改变 $x$，$\cat{C}(a, x)$ 也会在 $\Set$ 中变化。于是我们得到了从 $x$ 到 $\Set$ 的映射。

\begin{figure}[H]
  \centering
  \includegraphics[width=0.45\textwidth]{images/hom-set.jpg}
\end{figure}

\noindent
若我们想强调正在考虑的是第二个参数的映射，就使用记号 $\cat{C}(a, -)$，其中短横线作为参数占位符。

这种对象映射很容易扩展到态射映射。设范畴 $\cat{C}$ 中有任意两对象 $x$ 和 $y$ 之间的态射 $f$。
按上述定义，对象 $x$ 被映射到集合 $\cat{C}(a, x)$，对象 $y$ 被映射到 $\cat{C}(a, y)$。
若此映射是函子，则 $f$ 必须被映射到这两个集合之间的函数：$\cat{C}(a, x) \to \cat{C}(a, y)$

我们逐点定义这个函数，即分别针对每个参数。选择 $\cat{C}(a, x)$ 的任意元素作为参数——记作 $h$。
当端点匹配时态射可以复合。恰好 $h$ 的目标与 $f$ 的源匹配，因此它们的复合：
$$f \circ h \Colon a \to y$$
是从 $a$ 到 $y$ 的态射，因此属于 $\cat{C}(a, y)$。

\begin{figure}[H]
  \centering
  \includegraphics[width=0.45\textwidth]{images/hom-functor.jpg}
\end{figure}

\noindent
我们刚刚找到了从 $\cat{C}(a, x)$ 到 $\cat{C}(a, y)$ 的函数，它可以作为 $f$ 的像。
若没有混淆风险，我们将这个提升后的函数写作：$\cat{C}(a, f)$
及其作用于态射 $h$ 的形式：
$$\cat{C}(a, f) h = f \circ h$$
由于该构造适用于任何范畴，必然也适用于Haskell类型范畴。在Haskell中，这个Hom函子更广为人知的名字是\code{Reader}函子：

\src{snippet01}

\src{snippet02}
现在考虑如果我们固定同态集合的目标而非源会发生什么。换句话说，我们问的问题是：映射 $\cat{C}(-, a)$ 是否也是函子？
答案是肯定的，但它不是协变(covariant)而是逆变(contravariant)的。这是因为同样的端点匹配导致的是后复合(postcomposition)而非前复合(precomposition)。

我们已经在Haskell中见过这个逆变函子。我们称其为\code{Op}：

\src{snippet03}

\src{snippet04}
最后，若让两个对象都变化，我们就得到一个profunctor $\cat{C}(-, =)$，它在第一个参数逆变而在第二个参数协变
（为了强调两个参数独立变化，我们在第二个占位符使用双短横线）。
当我们讨论函子性时已经见过这个profunctor：

\src{snippet05}
重要结论是：这个观察适用于任何范畴：对象到同态集合的映射是函子性的。
由于逆变等价于从反范畴的映射，我们可以简洁地表示为：
$$C(-, =) \Colon \cat{C}^\mathit{op} \times \cat{C} \to \Set$$

\section{可表函子}

我们已经看到，对于范畴$\cat{C}$中任意对象$a$，
都能得到从$\cat{C}$到$\Set$的一个函子。这种到$\Set$
的保结构映射通常被称为\newterm{表示}（representation）。我们将$\cat{C}$的对象与态射表示为$\Set$中的集合与函数。

函子$\cat{C}(a, -)$本身有时被称为可表的。更一般地，若存在某个对象$a$使得函子$F$与Hom函子自然同构，则称$F$是\newterm{可表函子}（representable functor）。由于$\cat{C}(a, -)$是$\Set$-值函子，此类函子必然也是$\Set$-值的。

此前说过，我们常将同构的集合视为相同。更广泛地说，在范畴中我们将同构的对象视为相同，因为对象的结构完全由其通过态射与其他对象（及自身）的关系所定义。

以单子范畴$\cat{Mon}$为例，虽然最初用集合建模，但我们特别选取了保持单子结构的函数作为态射。因此当$\cat{Mon}$中两个对象同构（即存在可逆态射）时，它们具有完全相同的结构。若观察其底层集合与函数，会发现一个单子的单位元被映射到另一个单子的单位元，两个元素的乘积被映射为其像的乘积。

同样的推理适用于函子。函子间的自然变换扮演着态射角色，构成新的范畴。若存在可逆自然变换连接两个函子，则它们可被视为相同。

从这一视角分析可表函子的定义：若函子$F$可表，则需存在：
1. 范畴$\cat{C}$中的对象$a$；
2. 从$\cat{C}(a, -)$到$F$的自然变换$\alpha$；
3. 反向自然变换$\beta$；
4. 两者复合为恒等自然变换。

考察$\alpha$在对象$x$处的分量，它是$\Set$中的函数：
$$\alpha_x \Colon \cat{C}(a, x) \to F x$$
该变换的自然性条件表明：对任意$x$到$y$的态射$f$，下图交换：
$$F f \circ \alpha_x = \alpha_y \circ \cat{C}(a, f)$$
在Haskell中，自然变换对应多态函数：

\src{snippet06}
其中可选的\code{forall}量化符。由于参数化特性，自然性条件

\src{snippet07}
自动满足（这就是前文提到的"免费定理"之一），前提是左侧的\code{fmap}由函子$F$定义，右侧由读者函子定义。由于读者函子的\code{fmap}本质是函数预组合，我们可以更明确地表述：对$h \in \cat{C}(a, x)$应用自然性条件简化为：

\src{snippet08}
反向变换\code{beta}定义如下：

\src{snippet09}
它必须满足自然性条件且为\code{alpha}的逆：

\begin{snip}{text}
alpha . beta = id = beta . alpha
\end{snip}
后续将看到（Yoneda引理），只要$F a$非空，总存在从$\cat{C}(a, -)$到任意$\Set$-值函子的自然变换，但未必可逆。

以Haskell的列表函子和\code{Int}为例，以下自然变换成立：

\src{snippet10}
这里随意选择数字12构造单元素列表，再通过\code{h}进行\code{fmap}得到目标类型的列表。（实际上存在与整数列表数量相等的此类变换。）

自然性条件等价于\code{map}（列表版\code{fmap}）的可组合性：

\src{snippet11}
但尝试构造逆变换时，需要从任意类型\code{x}的列表转换为返回\code{x}的函数：

\src{snippet12}
或许考虑用\code{head}提取元素，但这对空列表无效。注意到无论选择何种类型\code{a}（替代\code{Int}）都无法解决此问题，故列表函子不可表。

回想之前将Haskell（自）函子类比为容器，类似地可将可表函子视为存储函数调用结果的备忘录容器。代表对象$a$（在$\cat{C}(a, -)$中）作为键类型，用于访问函数的表格化结果。我们称变换$\alpha$为\code{tabulate}，其逆$\beta$为\code{index}。以下是略微简化的\code{Representable}类定义：

\src{snippet13}
注意这里的代表类型$a$（即\code{Rep f}）属于\code{Representable}定义的一部分。星号表示\code{Rep f}是具体类型（而非类型构造器或其他更复杂的种类）。

无限列表（流）因永不为空而可表：

\src{snippet14}
可视作记忆了以\code{Integer}为参数的函数值。（严格来说应使用非负自然数，但为简化代码未作此限制。）

通过\code{tabulate}生成无限值流即可实现这类函数的记忆化。当然这得益于Haskell的惰性求值特性，值在需要时才被计算。使用\code{index}访问记忆化值：

\src{snippet15}
有趣的是，单一记忆化方案能覆盖具有任意返回类型的函数族。

反变函子的可表性定义类似，区别在于固定$\cat{C}(-, a)$的第二个参数。等价地可考虑从$\cat{C}^\mathit{op}$到$\Set$的函子，因为$\cat{C}^\mathit{op}(a, -)$等同于$\cat{C}(-, a)$。

可表性有个有趣的延伸：记得在笛卡尔闭范畴中Hom集可内部化为指数对象吗？Hom集$\cat{C}(a, x)$相当于$x^a$，对可表函子$F$可写作：$-^a = F$。

开个玩笑取对数两边得：$a = \mathbf{log}F$

当然这只是形式变换，但若了解对数性质则很有启发。特别地，基于积类型的函子可用和类型表示，而和类型函子通常不可表（例如列表函子）。

最后请注意，可表函子提供了同一事物的两种不同实现——一个是函数，一个是数据结构。它们具有完全相同的内容（通过相同键获取相同值），这正是前文所说的"相同性"。自然同构的函子在其内容上完全等价。然而两种表示常有不同实现方式，性能特征也可能差异显著。记忆化作为性能优化手段，可能大幅减少运行时间。能够生成同一底层计算的不同表示在实践中极具价值。令人惊讶的是，尽管范畴论本身不关注性能，却提供了大量探索实用替代实现的机会。

\section{挑战}

\begin{enumerate}
  \tightlist
  \item
        证明同态函子将范畴 \emph{C} 中的恒等态射映射到 $\Set$ 中相应的恒等函数。
  \item
        证明 \code{Maybe} 不是可表示函子。
  \item
        \code{Reader} 函子是可表示的吗？
  \item
        使用 \code{Stream} 的表示形式，对一个对其参数求平方的函数进行记忆化（memoization）。
  \item
        证明 \code{Stream} 的 \code{tabulate} 和 \code{index} 确实互为逆运算。（提示：使用数学归纳法。）
  \item
        考虑如下函子：

        \begin{snip}{haskell}
Pair a = Pair a a
\end{snip}
        它是可表示的。你能猜出代表它的类型吗？请实现 \code{tabulate} 和 \code{index}。
\end{enumerate}

\section{参考文献}

\begin{enumerate}
  \tightlist
  \item
        Catsters组合关于
        \urlref{https://www.youtube.com/watch?v=4QgjKUzyrhM}{可表函子}
        的视频。
\end{enumerate}